\documentclass[10pt,a4paper]{amsart}
\usepackage[margin=19mm]{geometry}
\usepackage{amsmath,amssymb,mathtools}
\usepackage[scr=rsfs,cal=euler]{mathalfa}
\usepackage{xfrac}
\usepackage[unicode,hidelinks,bookmarks=false]{hyperref}
\newcommand{\ie}{i.e.\ }
\newcommand{\cf}{cf.\ }
\newcommand{\resp}{respectively\ }
\newcommand{\Subsection}{\subsection}
\providecommand{\nobreakdash}{\nobreak\hbox{-}}
\setlength{\emergencystretch}{2em}
\multlinegap0pt
\usepackage{bm}
\usepackage{bbm}
\usepackage{dsfont}
\usepackage{xspace}
\usepackage{cancel}
\usepackage{mathtools}
\usepackage{amsthm}
\makeatletter
\@ifundefined{lemma}{\newtheorem{lemma}{Lemma}}{}
\makeatother
\title[A Non-compact Positivity-Preserving Scheme for Parabolic PDE]{A Non-compact Positivity-Preserving Scheme for Parabolic PDE via Conditional Expectation}
\author{Haoran Xu, Jie Ren, Xingye Yue}
\date{Source: arXiv:2601.10977v1 (CC BY 4.0)}
\begin{document}
\maketitle

\noindent\emph{Source and licence.} A Non-compact Positivity-Preserving Scheme for Parabolic PDE via Conditional Expectation. Version arXiv:2601.10977v1 (CC BY 4.0), distributed under the Creative Commons Attribution 4.0 International licence, as recorded in the arXiv licence metadata for this exact version and shown on the abstract page https://arxiv.org/abs/2601.10977v1. Full rights notice: licenses/arxiv-2601.10977-CC-BY-4.0.txt.\par\smallskip

\noindent\textit{Editorial scope.} Counted unit: the Lemma whose source label is \texttt{a} in the article (source0.tex lines 427--432, source label \texttt{a}), delivered together with the article's own complete original proof of it (source0.tex lines 434--514, one \texttt{proof} environment of 81 lines). No mathematics is written, repaired or re-derived: statement and proof are the article's own text line for line. Required context retained uncounted: eq:PDE-backward (lines 209--216); eq:weighted-scheme (lines 394--398); eq:omega (lines 400--411). Every definition, display and label of those lines is retained unchanged. Omitted from this package and not redistributed: 1--208, 217--393, 399--399, 412--426, 433--433, 515--1689. Nothing inside a retained excerpt is omitted or altered: every retained body is delivered exactly as the article's lines are written, so no omission receipt of any kind accompanies this package. Citations: the retained excerpts carry no citation command at all, so no bibliography entry of the article is transcribed and no reference is redistributed. Reference adaptation: none was needed and none was made; every reference inside the delivered unit resolves inside it, so no unresolved reference or citation is produced. Printed slips kept literal: no printed word, symbol, hypothesis or number of the counted statement or of its proof was altered. Layout and class: the article's own class is \texttt{article} with packages including geometry, setspace, amsmath, amssymb, amsthm, bm, mathrsfs, graphicx, hyperref, booktabs, titling; the wrapper is the lane's pinned \texttt{amsart} editorial wrapper at 10pt on A4, which restates the theorem-like environments the retained text needs and adds the article's own macro definitions that the retained text needs (none), with extra wrapper packages (bm, bbm, dsfont, xspace, cancel, mathtools). The delivered numbering is the unit's own. Assets: no retained span contains a figure environment, a graphics-inclusion command, a PDF-page inclusion, a table or a TikZ picture; no image file is copied into this package, the package declares no asset dependency and no image file is redistributed with it. Licence: version arXiv:2601.10977v1 of this article is distributed under the Creative Commons Attribution 4.0 International licence, as recorded in the arXiv licence metadata for this exact version and shown on the abstract page https://arxiv.org/abs/2601.10977v1. A full rights notice is delivered at licenses/arxiv-2601.10977-CC-BY-4.0.txt. Characters: every retained excerpt is pure ASCII, so the delivered engine, XeLaTeX with its default Latin Modern text font, reports no missing character.
	\begin{equation}\label{eq:PDE-backward}
		\begin{cases}
			\partial_t f + \dfrac{1}{2}\,\mathrm{Tr}\big( A A^\top D^2 f \big) + B\cdot\nabla f - r(x,y,t)\,f = 0,
			& (x,y)\in\Omega,\; t\in[0,T),\\[2pt]
			f(x,y,T) = \varphi(x,y), & (x,y)\in\Omega,\\[1pt]
			f\big|_{\partial\Omega} = f_{\partial}(x,y,t), & t\in[0,T),
		\end{cases}
	\end{equation}

	\begin{equation}\label{eq:weighted-scheme}
		\hat f_{h,i,j}^n
		= \sum_{k=1}^4 \frac{\omega_k}{\,1+r_{i,j}^n\hat\tau_k\,}\,
		f(X_{\tau_k}^{h,k}, Y_{\tau_k}^{h,k}, \tau_k),
	\end{equation}

	\begin{equation}\label{eq:omega}
		\begin{aligned}
			\omega_1 &= \frac{\sqrt{\hat\tau_2 \hat\tau_3 \hat\tau_4}}
			{(\sqrt{\hat\tau_1}+\sqrt{\hat\tau_3})(\sqrt{\hat\tau_1 \hat\tau_3}+\sqrt{\hat\tau_2 \hat\tau_4})},\\
			\omega_2 &= \frac{\sqrt{\hat\tau_1 \hat\tau_3 \hat\tau_4}}
			{(\sqrt{\hat\tau_2}+\sqrt{\hat\tau_4})(\sqrt{\hat\tau_1 \hat\tau_3}+\sqrt{\hat\tau_2 \hat\tau_4})},\\
			\omega_3 &= \frac{\sqrt{\hat\tau_1 \hat\tau_2 \hat\tau_4}}
			{(\sqrt{\hat\tau_1}+\sqrt{\hat\tau_3})(\sqrt{\hat\tau_1 \hat\tau_3}+\sqrt{\hat\tau_2 \hat\tau_4})},\\
			\omega_4 &= \frac{\sqrt{\hat\tau_1 \hat\tau_2 \hat\tau_3}}
			{(\sqrt{\hat\tau_2}+\sqrt{\hat\tau_4})(\sqrt{\hat\tau_1 \hat\tau_3}+\sqrt{\hat\tau_2 \hat\tau_4})}.
		\end{aligned}
	\end{equation}

	\begin{lemma}\label{a}
		Assume that the exact solution $f$ is sufficiently smooth. Let $\hat f_{h,i,j}^n$ be the numerical approximation obtained from the weighted scheme \eqref{eq:weighted-scheme}. Then there exists a constant $C>0$ independent of $\Delta t$ such that
		\[
		\bigl|\hat f_{h,i,j}^n - f(x_i, y_j, t_n)\bigr| \;\le\; C\,\Delta t^{3/2}.
		\]
	\end{lemma}

	\begin{proof}
		We begin by defining the quadrature error term $I$ as
		\[
		\begin{aligned}
			I &= \hat f_{h,i,j}^n - f(x_i, y_j, t_n) \\
			&= \sum_{k=1}^4 \frac{\omega_k}{\,1 + r_{i,j}^n\hat{\tau}_k\,}\,
			f(X_{\tau_k}^{h,k}, Y_{\tau_k}^{h,k}, \tau_k)
			- f(x_i, y_j, t_n)  \\
			&= \Bigl[ A_1\,\partial_t f + A_2\,\partial_{xx}^2 f + A_3\,\partial_{xy}^2 f
			+ A_4\,\partial_{yy}^2 f + A_5\,\partial_x f + A_6\,\partial_y f + A_7\,f \Bigr]_{(x_i,y_j)}^{t_n}
			+ O(\Delta t^{3/2}),
		\end{aligned}
		\]
		where the coefficients $A_1,\dots,A_7$ are explicit algebraic combinations of the probabilities $\{\omega_k\}$, the time increments $\{\hat\tau_k\}$, and the local coefficients:
		\[
		\begin{aligned}
			A_1 &= \omega_1\,\hat{\tau}_1 + \omega_2\,\hat{\tau}_2
			+ \omega_3\,\hat{\tau}_3 + \omega_4\,\hat{\tau}_4,\\
			A_2 &= \tfrac{(\sigma_{1,i,j}^n)^2}{2}\,A_1 + \tfrac{\rho_{i,j}^n(\sigma_{1,i,j}^n)^2}{2}\,
			(\omega_1\,\hat{\tau}_1 - \omega_2\,\hat{\tau}_2 + \omega_3\,\hat{\tau}_3 - \omega_4\,\hat{\tau}_4),\\
			A_3 &= \rho_{i,j}^n\sigma_{1,i,j}^n\sigma_{2,i,j}^n\,A_1
			+ \sigma_{1,i,j}^n\sigma_{2,i,j}^n\,
			(\omega_1\,\hat{\tau}_1 - \omega_2\,\hat{\tau}_2 + \omega_3\,\hat{\tau}_3 - \omega_4\,\hat{\tau}_4),\\
			A_4 &= \tfrac{(\sigma_{2,i,j}^n)^2}{2}\,A_1 + \tfrac{\rho_{i,j}^n(\sigma_{2,i,j}^n)^2}{2}\,
			(\omega_1\,\hat{\tau}_1 - \omega_2\,\hat{\tau}_2 + \omega_3\,\hat{\tau}_3 - \omega_4\,\hat{\tau}_4),\\
			A_5 &= b_{1,i,j}^n\,A_1 + \omega_1\,\alpha_{i,j}^n\,\sigma_{1,i,j}^n\,\sqrt{\hat\tau_1}
			- \omega_2\,\beta_{i,j}^n\,\sigma_{1,i,j}^n\,\sqrt{\hat\tau_2}\\
			&\quad - \omega_3\,\alpha_{i,j}^n\,\sigma_{1,i,j}^n\,\sqrt{\hat\tau_3}
			+ \omega_4\,\beta_{i,j}^n\,\sigma_{1,i,j}^n\,\sqrt{\hat\tau_4},\\
			A_6 &= b_{2,i,j}^n\,A_1 + \omega_1\,\alpha_{i,j}^n\,\sigma_{2,i,j}^n\,\sqrt{\hat\tau_1}
			+ \omega_2\,\beta_{i,j}^n\,\sigma_{2,i,j}^n\,\sqrt{\hat\tau_2}\\
			&\quad - \omega_3\,\alpha_{i,j}^n\,\sigma_{2,i,j}^n\,\sqrt{\hat\tau_3}
			- \omega_4\,\beta_{i,j}^n\,\sigma_{2,i,j}^n\,\sqrt{\hat\tau_4},\\
			A_7 &= -\,r_{i,j}^n\,A_1.
		\end{aligned}
		\]
		
		From \eqref{eq:omega}, it follows that the probabilities $\{\omega_k\}$ satisfy the moment-matching conditions:
		\[
		\begin{aligned}
			\omega_1 + \omega_2 + \omega_3 + \omega_4 &= 1,\\
			\omega_1\,\hat{\tau}_1 - \omega_2\,\hat{\tau}_2
			+ \omega_3\,\hat{\tau}_3 - \omega_4\,\hat{\tau}_4 &= 0,\\
			\omega_1\,\alpha_{i,j}^n\,\sigma_{1,i,j}^n\,\sqrt{\hat\tau_1}
			- \omega_2\,\beta_{i,j}^n\,\sigma_{1,i,j}^n\,\sqrt{\hat\tau_2}
			- \omega_3\,\alpha_{i,j}^n\,\sigma_{1,i,j}^n\,\sqrt{\hat\tau_3}
			+ \omega_4\,\beta_{i,j}^n\,\sigma_{1,i,j}^n\,\sqrt{\hat\tau_4} &= 0,\\
			\omega_1\,\alpha_{i,j}^n\,\sigma_{2,i,j}^n\,\sqrt{\hat\tau_1}
			+ \omega_2\,\beta_{i,j}^n\,\sigma_{2,i,j}^n\,\sqrt{\hat\tau_2}
			- \omega_3\,\alpha_{i,j}^n\,\sigma_{2,i,j}^n\,\sqrt{\hat\tau_3}
			- \omega_4\,\beta_{i,j}^n\,\sigma_{2,i,j}^n\,\sqrt{\hat\tau_4} &= 0.
		\end{aligned}
		\]
		
		Substituting these identities into the expressions for $A_2$--$A_7$, we conclude that
		\[
		\begin{aligned}
			A_2 &= \tfrac{(\sigma_{1,i,j}^n)^2}{2}\,A_1,\qquad 
			A_3 = \rho_{i,j}^n\,\sigma_{1,i,j}^n\,\sigma_{2,i,j}^n\,A_1,\qquad
			A_4 = \tfrac{(\sigma_{2,i,j}^n)^2}{2}\,A_1,\\
			A_5 &= b_{1,i,j}^n\,A_1,\qquad 
			A_6 = b_{2,i,j}^n\,A_1,\qquad 
			A_7 = -\,r_{i,j}^n\,A_1.
		\end{aligned}
		\]
		
		Therefore, the error term simplifies to
		\[
		I = A_1\Bigl[\partial_t f + \tfrac{\sigma_1^2}{2}\,\partial_{xx}^2 f
		+ \rho\,\sigma_1\,\sigma_2\,\partial_{xy}^2 f 
		+ \tfrac{\sigma_2^2}{2}\,\partial_{yy}^2 f
		+ b_1\,\partial_x f + b_2\,\partial_y f - r\,f\Bigr]_{(x_i,y_j)}^{t_n}
		+ O(\Delta t^{3/2}).
		\]
		
		Since $f$ satisfies the original PDE \eqref{eq:PDE-backward}, the bracketed expression vanishes. Hence
		\[
		I = O(\Delta t^{3/2}),
		\]
		which completes the proof.
	\end{proof}

\end{document}
